\title{Direct Preference Optimization: Your Language Model is Secretly a Reward Model}

\begin{abstract}
Large-scale unsupervised language models (LMs) acquire extensive world knowledge and demonstrate certain reasoning abilities; however, effectively directing their behavior remains challenging because of the unsupervised approach underlying their training. To improve such steerability, current strategies typically rely on collecting human judgments that compare different model outputs and subsequently adapting the LM through preference alignment, often using reinforcement learning from human feedback (RLHF). Nonetheless, RLHF involves a multi-step and frequently unstable process: first, a reward model is trained to represent human preferences, and then the original LM is optimized via reinforcement learning to maximize this inferred reward while maintaining proximity to the initial model parameters. In this work, we propose a novel reward model parameterization within the RLHF framework, which makes it possible to analytically derive the optimal policy, thereby reframing the conventional RLHF problem into one solvable using a straightforward classification objective. Our method, termed \textit{Direct Preference Optimization} (DPO), offers stability, strong empirical performance, and computational efficiency by obviating the need for sampling from the LM during fine-tuning or for extensive hyperparameter tuning. Through empirical evaluation, we demonstrate that DPO fine-tunes LMs to match or surpass the effectiveness of existing preference alignment methods. In particular, DPO fine-tuning provides superior sentiment control compared to PPO-based RLHF, while matching or exceeding generation quality for tasks like summarization and single-turn dialogue, all with a simpler and more accessible implementation and training process.
\end{abstract}

\section{Introduction}
Large-scale unsupervised language models (LMs) trained on extensive corpora exhibit remarkable emergent capabilities~\citep{chowdhery2022palm, brown2020language, touvron2023llama,bubeck2023sparks}. Nevertheless, these models are exposed during training to data originating from humans who possess a diverse array of intentions, objectives, and expertise. Not all of these attributes are suitable for imitation; for instance, while an AI code assistant should ideally \textit{recognize} prevalent programming errors in order to provide corrections, its code generation ought to reflect the higher-quality coding proficiency—possibly infrequently represented—within the dataset. In a related vein, a language model might benefit from being \textit{cognizant} of a widespread misconception affecting half of individuals, but we would strongly prefer that it not reproduce this misconception as a factual claim in half the relevant queries it receives. Accordingly, it is imperative to distinguish and promote the model's \emph{desired responses and behaviors} from the broad array of its \textit{knowledge and skills} to create AI systems that are both reliable and manageable~\citep{ouyang2022training}. Although conventional approaches typically guide LMs to align with human values through reinforcement learning (RL), we demonstrate that the RL-based preference optimization objective in standard methods can, in fact, be implemented precisely using a straightforward binary cross-entropy loss, thus streamlining the entire preference learning process.

In summary, current techniques shape the target behaviors of a language model using curated datasets composed of human preferences, which encode traits considered beneficial and safe. This phase of learning from preference annotations typically succeeds an initial period of large-scale unsupervised pre-training on massive text datasets. While the simplest preference learning strategy involves supervised fine-tuning using high-quality human response examples, the most prominent and effective paradigm remains reinforcement learning from human (or AI) feedback (RLHF/RLAIF; \citep{christiano2017deep,bai2022constitutional}). RLHF methods construct a reward model from human preference datasets and leverage RL to train the language model to produce responses with high associated rewards, while constraining deviation from the original model parameters. Despite the substantial conversational and coding competence enabled by RLHF, this framework is much more intricate than supervised fine-tuning. It necessitates training several language models and frequent sampling from the policy during training iterations, thereby resulting in considerable computational overhead.

This work demonstrates a method for optimizing language models to reflect human preferences without the need for explicit reward modeling or reinforcement learning techniques. We introduce \textit{Direct Preference Optimization (DPO)}, a straightforward algorithm that implicitly achieves the same goal as standard RLHF approaches—namely, reward maximization under a KL-divergence constraint—while offering ease of implementation and training. At a high level, the DPO algorithm updates the model by increasing the ratio of the log probability assigned to preferred over dispreferred outputs, leveraging an adaptive, per-instance weighting scheme. This scheme counteracts the model degradation observed when applying a simple probability ratio objective. DPO, like comparable algorithms, is grounded in a theoretical preference framework such as the Bradley-Terry model \cite{bradley1952rankanalysis}, which assesses the concordance between a reward function and human preference observations. However, in contrast to conventional approaches that utilize the preference model to construct a loss for training a reward model—subsequently used to guide policy training—DPO employs a variable transformation so that the preference loss directly depends on the policy itself. Given a collection of annotated human preferences between model outputs, DPO is capable of optimizing the policy via a binary cross-entropy loss, thereby yielding an optimal policy corresponding to an implicit reward function tailored to the preference data.

The principal innovation of this paper is Direct Preference Optimization (DPO), an intuitive algorithm that dispenses with reinforcement learning for training language models from preference feedback. Empirically, we show that DPO matches or surpasses the performance of prior techniques, including those based on PPO-style RLHF, for various preference-driven applications such as sentiment modification, summarization, and dialogue—using models containing up to 6B parameters.

Self-supervised language models, particularly as their scale increases, exhibit an emergent ability to solve certain tasks either without any explicit supervision (zero-shot) \citep{radford2019language} or with only minimal prompting (few-shot) \citep{gpt3,megatron,chowdhery2022palm}. Nevertheless, their efficacy on downstream tasks as well as their alignment with user intentions is often greatly enhanced through fine-tuning on corpora composed of instructions paired with human-authored completions \citep{mishra-etal-2022-cross,sanh2022multitask,chung2022scaling,thoppilan2022lamda}. This procedure, known as `instruction-tuning,' equips LLMs with improved generalization capabilities to handle instructions beyond those seen during fine-tuning, resulting in more broadly usable systems \citep{chung2022scaling}. Despite the advances achieved through instruction tuning, it is frequently more feasible to obtain \textit{relative} assessments of completion quality from humans than to assemble large collections of high-quality demonstrations. Consequently, recent approaches have focused on further fine-tuning LLMs using datasets grounded in human preferences, leading to advancements in tasks such as translation \citep{kreutzer-etal-2018-reliability}, summarization \citep{stiennon2022learning,ziegler2020finetuning}, narrative generation \citep{ziegler2020finetuning}, and instruction compliance \citep{ouyang2022training,ramamurthy2023is}. The standard practice involves first learning a neural network-based reward model designed to reflect observed human preference data—often relying on frameworks like the Bradley-Terry model \citep{bradley1952rankanalysis}—followed by optimizing the language model to maximize the learned reward via reinforcement learning algorithms such as REINFORCE \citep{williams1992reinforce}, proximal policy optimization (PPO; \cite{schulman2017proximal}), or related variants \citep{ramamurthy2023is}. A related stream of research utilizes LLMs, themselves instruction-tuned with feedback from humans, to autonomously generate additional preference-labeled data regarding properties like safety or non-harmfulness \citep{bai2022constitutional}, where weak human oversight takes the form of textual rubrics for the LLM’s outputs. Collectively, these approaches bring together insights from the application of reinforcement learning to train language models for complex objectives~\citep{Ranzato2015SequenceLT,paulus2018a,wu2018learning} and from general methodologies for preference-based learning from human feedback \citep{christiano2017deep,kupcsik2018learning}. Nonetheless, while relative human preferences offer practical advantages, the reinforcement learning-based fine-tuning of large language models poses considerable practical hurdles; the present work introduces a theoretically grounded method for optimizing language models with respect to relative preferences, bypassing the need for reinforcement learning.

Independent of linguistic applications, research has investigated policy learning from preferences within both the bandit and reinforcement learning paradigms, proposing diverse methodologies. A notable example is contextual dueling bandits (CDB; \cite{yue2012karmed,dudik2015contextual}), wherein preferences or rankings over actions substitute for explicit rewards. Lacking access to absolute reward signals, CDB theory replaces the notion of an optimal policy with that of a \textit{von Neumann winner}, defined as a policy that defeats \textit{every} other policy with a probability of at least 50\% in expectation \citep{dudik2015contextual}. A key distinction, however, is that CDBs acquire preference labels sequentially in an online manner, in contrast to human preference learning, which typically involves optimizing from a finite, static set of offline, preference-annotated action pairs \citep{yan2022human}. Analogously, \textit{preference-based RL} (PbRL) operates on binary preferences informed by an \textit{unknown} scoring function rather than by explicit reward feedback \citep{BusaFekete2014,ruiz2023dueling}. Several algorithms have been proposed for PbRL, some of which support the reuse of off-policy preference data; however, these strategies commonly involve an explicit phase of learning the hidden scoring or reward model, followed by optimization with respect to this estimate \citep{jain2013learning,BusaFekete2014,christiano2017deep,sadigh2017active,kupcsik2018learning}. By contrast, our method involves a unified policy learning procedure that aims to directly optimize the policy to comply with observed preferences.

\section{Preliminaries}\label{section:prelims}

Here, we outline the RLHF framework as presented in \citeauthor{ziegler2020finetuning}, subsequently followed in works such as \citep{stiennon2022learning, bai2022training, ouyang2022training}. This pipeline comprises three principal stages: (1) supervised fine-tuning (SFT); (2) preference-based sampling and reward modeling; and (3) reinforcement learning optimization.

\textbf{SFT}: The RLHF process starts by employing supervised learning to further tune a pre-trained language model using a curated dataset specific to the downstream application (such as dialogue generation or summarization), yielding the SFT model, denoted $\pi^\text{SFT}$.

\textbf{Reward Modelling Phase}: Next, the SFT model is conditioned on prompts $x$ to sample answer pairs $(y_1, y_2)\sim \pi^\text{SFT}(y \mid x)$. Human annotators are then presented with these pairs and asked to indicate their preference, written as $y_w\succ y_l \mid x$, where $y_w$ denotes the response deemed superior to $y_l$ given prompt $x$. These preferences are presumed to originate from a latent, unknown reward model $r^*(y, x)$. Several probabilistic models can be employed to model such preferences; the Bradley-Terry (BT) model \cite{bradley1952rankanalysis} is a frequently utilized approach, though broader models such as Plackett-Luce ranking models \citep{plackett1975analysis, luce2012individual} are compatible if multiple ranked outputs are provided. According to the BT model, the human preference probability distribution $p^*$ is formulated as follows:

Given a fixed dataset of pairwise comparisons $\mathcal{D}=\bigl\{x^{(i)}, y_w^{(i)}, y_l^{(i)}\bigr\}_{i=1}^N$ sampled according to $p^*$, we define a reward function $r_{\phi}(x, y)$, parameterized by $\phi$, which can be fitted using maximum likelihood estimation. Treating the task as binary classification, the loss function corresponding to the negative log-likelihood is:

\begin{equation}\label{eq:reward_model}
    \mathcal{L}_R(r_{\phi}, \mathcal{D}) = -\mathbb{E}_{(x, y_w, y_l)\sim \mathcal{D}}\left[\log \sigma\bigl(r_{\phi}(x, y_w)- r_{\phi}(x, y_l)\bigr)\right]
\end{equation}

where $\sigma$ denotes the logistic sigmoid. In the scenario of large language models, $r_{\phi}(x, y)$ typically builds upon the SFT model $\pi^\text{SFT}(y \mid x)$ and incorporates an additional linear output head after the final transformer block to yield a scalar reward, as described in \cite{ziegler2020finetuning}. For the sake of stability, normalization techniques are frequently employed to enforce $\mathbb{E}_{x,y\sim \mathcal{D}}\left[r_\phi(x, y)\right] = 0$ for all $x$, thus reducing variance in reward predictions.

\textbf{RL Fine-Tuning Phase}: In the subsequent reinforcement learning stage, the policy is updated using feedback from the trained reward model. Consistent with earlier research~\citep{jaques2017sequence, jaques2020human}, policy optimization is given by

\begin{equation}\label{eq:RL}
\max_{\pi_{\theta}}  \mathbb{E}_{x\sim \mathcal{D}, y\sim \pi_{\theta}(y \mid x)}\left[r_{\phi}(x, y)\right] - \beta\mathbb{D}_{\textrm{KL}}\bigl[\pi_{\theta}(y\mid x)\mid \mid \pi_\text{ref}(y\mid x)\bigr],
\end{equation}

where $\beta$ modulates how closely $\pi_{\theta}$ tracks the reference policy $\pi_\text{ref}$ (i.e., the SFT baseline $\pi^\text{SFT}$). It is common practice to initialize $\pi_{\theta}$ with $\pi^\text{SFT}$. The regularization term is critical as it ensures that the optimized policy remains close to the data regime where the reward model is valid, and safeguards against loss of output diversity and mode collapse to singular, high-scoring responses. Since the problem of text generation involves discrete outputs, the objective is inherently non-differentiable and thus is usually tackled using reinforcement learning methods. A prevalent strategy \citep{ziegler2020finetuning, stiennon2022learning, bai2022training, ouyang2022training} is to define the composite reward as ${r(x, y) = r_{\phi}(x, y) -\beta (\log \pi_{\theta}(y\mid x) - \log \pi_\text{ref}(y\mid x))}$ and to optimize this function utilizing PPO \cite{schulman2017proximal}.

\section{Direct Preference Optimization}\label{sec:DPO}

In light of the significant computational hurdles associated with reinforcement learning-based tuning of large-scale language models, we pursue a streamlined approach for optimizing policies via preferences alone. In contrast with traditional RLHF strategies, where a reward model is first fitted and then maximized through RL, our methodology chooses a reward parameterization that allows for exact, closed-form policy extraction, eliminating the need for iterative RL training. The central idea is to use an explicit mathematical correspondence between reward structures and optimal policies, thereby recasting loss functions originally defined for rewards as equivalent losses for policy parameters. This variable transformation approach lets us train policies directly using pairwise preference data without training an independent reward model, but still preserves compatibility with human preference frameworks such as the Bradley-Terry model. Accordingly, the policy model subsumes both the language modeling and the implicit reward assignment within a single network.

\textbf{Derivation of the DPO objective.} We begin with the reinforcement learning objective given by Eq.~\ref{eq:RL}, following previous literature, for an arbitrary reward function $r$. As established in prior studies~\citep{peters2007reinforcement, peng2019advantage, korbak2022reinforcement, go2023aligning}, it can be readily demonstrated that the KL-regularized reward optimization problem defined in Eq.~\ref{eq:RL} yields an optimal solution of the form:

\begin{equation}\label{eq:op_policy}
    \pi_r(y\mid x) = \frac{1}{Z(x)}\pi_\text{ref}(y\mid x)\exp\left(\frac{1}{\beta}r(x, y)\right),
\end{equation}

where the normalization constant is $Z(x) =\sum_{y}\pi_\text{ref}(y\mid x)\exp\left(\frac{1}{\beta}r(x, y)\right)$. The complete derivation is included in Appendix \ref{app:derivation1}. Computing $Z(x)$ directly, even when utilizing an MLE-based reward model $r_\phi$ in place of the true reward $r^*$, remains computationally demanding~\citep{korbak2022reinforcement, go2023aligning}, limiting the practical application of this formulation. To address this, Eq.~\ref{eq:op_policy} can be rearranged to express the reward as a function of the associated optimal policy $\pi_r$, the reference policy $\pi_\text{ref}$, and the partition function $Z(\cdot)$. In particular, by taking the logarithm of both sides of Eq.~\ref{eq:op_policy} and isolating terms, we find:

\begin{equation}\label{eq:main_eq}
    r(x,y) =\beta \log \frac{\pi_r(y\mid x)}{\pi_\text{ref}(y\mid x)} + \beta \log Z(x).
\end{equation}

This alternative parameterization can be similarly applied to the true reward $r^*$ and its optimal policy $\pi^*$. Notably, the Bradley-Terry framework considers only the reward difference between pairs of candidate outputs, that is, ${p^*(y_1 \succ y_2 \mid x) = \sigma(r^*(x, y_1) - r^*(x, y_2))}$. Substituting the relation from Eq.~\ref{eq:main_eq} for $r^*(x, y)$ into the preference expression Eq.~\ref{eq:bradley-terry}, the partition function cancels out, allowing the human preference probability to be expressed solely using the optimal policy $\pi^*$ and reference policy $\pi_\text{ref}$. Thus, the optimal policy $\pi^*$ obtained by RLHF under the Bradley-Terry paradigm fulfills the following preference likelihood:

\begin{equation}\label{eq:objective}
    p^*(y_1\succ y_2 \mid x)=\frac{1}{1 + \exp\left(\beta \log \frac{\pi^*(y_2\mid x)}{\pi_\text{ref}(y_2\mid x)} - \beta \log \frac{\pi^*(y_1\mid x)}{\pi_\text{ref}(y_1\mid x)}\right)}
\end{equation}

Full details of this derivation can be found in Appendix~\ref{app:derivation2}. Although this expression uses the Bradley-Terry model, analogous constructions for more general preference models such as Plackett-Luce~\citep{plackett1975analysis, luce2012individual} are provided in Appendix~\ref{app:plackett_luce_models}.

Having reformulated the probability of observed human preferences in terms of the optimal policy rather than the underlying reward, we are now able to specify a maximum likelihood estimation objective for a policy parameterized by $\pi_\theta$. Mirroring the approach employed for reward model training (cf. Eq.~\ref{eq:reward_model}), the policy-based objective is as follows:

\begin{equation}\label{eq:optimum_model}
    \mathcal{L}_\text{DPO}(\pi_{\theta}; \pi_\text{ref}) = -\mathbb{E}_{(x, y_w, y_l)\sim \mathcal{D}}\left[\log \sigma \left(\beta \log \frac{\pi_{\theta}(y_w\mid x)}{\pi_\text{ref}(y_w\mid x

In this formulation, we estimate an implicit reward through an alternative parameterization, where the optimal policy coincides with $\pi_\theta$. Furthermore, because this approach amounts to fitting a reparameterized Bradley-Terry model, it inherits theoretical guarantees, such as consistency under appropriate assumptions on the preference data distribution \cite{bong2022generalized}. We elaborate on these theoretical aspects of DPO and its relationship to other methods in Section~\ref{sec:theory}.

\textbf{Mechanistic perspective on the DPO update.} To better understand the mechanics underlying DPO, we examine the gradient of its loss function, $\mathcal{L}_\text{DPO}$. The gradient with respect to the model parameters $\theta$ is given by:

\begin{multline*}\label{eq:gradient}
    \nabla_\theta \mathcal{L}_\text{DPO}(\pi_\theta;\pi_\text{ref}) = \\ -\beta\mathbb{E}_{(x, y_w, y_l) \sim \mathcal{D}} \bigg[\underbrace{\sigma(\hat{r}_\theta(x, y_l) - \hat{r}_\theta (x, y_w))}_\text{greater influence when reward prediction errs}\bigg[\underbrace{\nabla_\theta\log \pi(y_w \mid x)}_\text{boost probability of $y_w$} - \underbrace{\nabla_\theta\log\pi(y_l \mid x)}_\text{reduce probability of $y_l$}\bigg]\bigg],
\end{multline*}

Here, $\hat{r}_\theta(x, y) = \beta \log \frac{\pi_\theta(y \mid x)}{\pi_\text{ref}(y \mid x)}$ serves as the reward implicitly assigned by the language model $\pi_\theta$ relative to the reference $\pi_\text{ref}$ (see Section~\ref{sec:theory} for further discussion). Conceptually, the gradient direction increases the likelihood of favorable responses $y_w$, while reducing the likelihood for less-preferred outputs $y_l$. Crucially, each training example is weighted in proportion to the extent that the implicit reward $\hat{r}_\theta$ overestimates the less-preferred completion, modulated by $\beta$. In other words, the influence is larger when the current policy misorders the completions, capturing both the model error and the regularization imposed by the KL constraint. Empirically, our results underscore the necessity of this weighting: omitting the weighting factor can lead the language model to undesirable behavior (see Appendix Table~\ref{tab:unlikelihood_generations}).

\textbf{DPO overview.} The typical DPO workflow involves the following steps: 1) For each prompt $x$, generate samples $y_1, y_2 \sim \pi_\text{ref}(\cdot \mid x)$, and assign human preference labels to assemble an offline preference dataset $\mathcal{D} = \{x^{(i)}, y_w^{(i)}, y_l^{(i)}\}_{i=1}^N$; 2) Update the language model $\pi_\theta$ by minimizing $\mathcal{L}_\text{DPO}$ with respect to the fixed $\pi_\text{ref}$, preference data $\mathcal{D}$, and selected value of $\beta$.  
In actual scenarios, it is often preferable to utilize existing publicly available preference datasets rather than generate new outputs and collect human preference annotations. Since such datasets are usually derived from $\pi^\text{SFT}$, we set $\pi_\text{ref} = \pi^\text{SFT}$ when it is accessible. Conversely, if $\pi^\text{SFT}$ is absent, we construct $\pi_\text{ref}$ by maximizing the likelihood over preferred completions ${(x, y_w)}$, specifically, by setting ${\pi_\text{ref} = \argmax_{\pi}\mathbb{E}_{x, y_w \sim \mathcal{D}}\left[\log \pi(y_w \mid x)\right]}$. This approach aims to reduce discrepancies between the unattainable true reference distribution and the $\pi_\text{ref}$ adopted in DPO. Comprehensive implementation details and hyperparameter choices are discussed in Appendix~\ref{app:implementation}.

\section{Theoretical Analysis of DPO} This section aims to further elucidate the DPO framework, present a theoretical foundation, and connect the benefits of DPO to limitations observed in actor-critic approaches commonly employed in RLHF settings (for instance, PPO~\cite{schulman2017proximal}).

\label{sec:theory}

\subsection{Your Language Model Is Secretly a Reward Model} DPO is distinct in that it achieves both implicit reward modeling and policy optimization by leveraging a single maximum likelihood loss, bypassing the need to learn an explicit reward model or perform RL to adjust the policy. Importantly, the objective function (Eq. \ref{eq:main_eq}) can be reformulated as a Bradley-Terry model with rewards defined as $r^*(x, y) = \beta \log\frac{\pi^*_\theta(y \mid x)}{\pi_\text{ref}(y \mid x)}$. Here, training the parametric model $\pi_\theta$ is equivalent, under a variable substitution, to fitting the reward model objective in Eq. \ref{eq:reward_model}. In this part, we construct the theoretical underpinnings of this reparameterization, demonstrate that it does not reduce the space of representable reward functions, and argue that it enables precise recovery of the optimal policy. We commence by formalizing an equivalence relation among reward functions.

\begin{definition}
Two reward functions $r(x, y)$ and $r'(x, y)$ are defined as equivalent if and only if there exists a function $f$ such that $r(x, y)-r'(x, y) = f(x)$.     
\end{definition}

It follows readily that this definition establishes an equivalence relation, thereby dividing the set of reward functions into equivalence classes. The following lemmas summarize key properties:

\begin{lemma}\label{lemma:same_prefrence} Within the Plackett-Luce framework, including the specific case of Bradley-Terry preferences, any two reward functions from the same equivalence class yield identical preference distributions.
\end{lemma}

\begin{lemma}\label{lemma:same_policy}
    If two reward functions are from the same equivalence class, they result in the same optimal policy in the corresponding constrained RL formulation.
\end{lemma}

Proofs for these statements are uncomplicated and can be found in Appendix \ref{app:lemma1}. The first lemma captures a well-known identifiability issue inherent in the Plackett-Luce family \cite{plackett1975analysis}. To address this ambiguity, it is customary to impose identifiability constraints in order to secure uniqueness of the maximum likelihood estimates, as discussed in Eq.~\ref{eq:reward_model} and related literature \cite{bong2022generalized}. The second lemma demonstrates that all members of a given reward class produce the same optimal policy. Thus, for practical purposes, it is sufficient to recover any representative from the equivalence class that produces the optimal policy. We formally establish the following result in Appendix~\ref{app:thm1}:

\begin{theorem}\label{thm:main}
    Assuming mild conditions, every reward class consistent with Plackett-Luce or, specifically, the Bradley-Terry models, admits a representation of the form ${r(x, y) = \beta \log \frac{\pi(y\mid x)}{\pi_\text{ref}(y\mid x)}}$ for some model $\pi(y\mid x)$ and specified reference model $\pi_\text{ref}(y \mid x)$.
\end{theorem}

\begin{sproof}
    Let $r(x, y)$ be an arbitrary reward function, inducing an associated optimal policy $\pi_r(y \mid x)$ via Eq.~\ref{eq:op_policy}. We now argue that some reward function in the equivalence class of $r$ can be written using the specified reparameterization. Consider the projection defined by
\begin{equation}
    f(r; \pi_\text{ref}, \beta)(x, y) = r(x, y) - \beta\log\sum_{y}\pi_\text{ref}(y\mid x)\exp\left(\frac{1}{\beta}r(x, y)\right)
\end{equation}
This operator, $f$, effectively centers the reward function using the log-partition function of $\pi_r$. Since this normalization depends solely on $x$, $f(r; \pi_\text{ref}, \beta)(x, y)$ remains within the same equivalence class as $r(x, y)$. Substituting the right-hand side of Eq.~\ref{eq:main_eq} (which characterizes all reward functions), we see that $f(r; \pi_\text{ref}, \beta)(x, y) = \beta \log \frac{\pi_r(y\mid x)}{\pi_\text{ref}(y\mid x)}$. Thus, this projection selects a representative within the class of $r$ exhibiting the targeted structural form, confirming the full generality of the proposed reparameterized reward model.
\end{sproof}

Theorem~\ref{thm:main} can also be interpreted as precisely identifying which representative of the reward function equivalence class is chosen by the DPO reparameterization. Specifically, it selects the reward function satisfying the condition:

\begin{equation}\label{eq:lag_p}
     \sum_{y}\underbrace{\pi_\text{ref}(y\mid x)\exp\left(\frac{1}{\beta}r(x, y)\right)}_{=\pi(y\mid x)\text{, using Thm.~\ref{thm:main} reparam.}} = 1,
\end{equation}

that is, $\pi(y\mid x)$ forms a legitimate probability distribution (all probabilities are nonnegative and their total sum is 1). Notably, by referencing Eq.~\ref{eq:op_policy}, we observe that Eq.~\ref{eq:lag_p} characterizes the partition function of the optimal policy determined by the reward function $r(x, y)$. The fundamental contribution of the DPO algorithm lies in its capacity to introduce specific constraints to the under-specified Plackett-Luce (and more particularly, Bradley-Terry) preference models. These constraints do not diminish the richness of the representable reward function classes, but instead guarantee that the optimal policy described in Eq.~\ref{eq:op_policy} remains explicitly solvable for every prompt $x$.

\subsection{Instability of Actor-Critic Algorithms} Our framework further enables the identification and analysis of sources of instability in commonly used actor-critic methods for RLHF, including PPO. Adhering to the standard RLHF procedure, we concentrate on the RL fine-tuning phase described in Section \ref{section:prelims}. The relationship between this process and the control as inference approach to constrained RL, as discussed in \cite{levine2018reinforcement}, and problem \ref{eq:RL}, can be made explicit. Assuming a parameterized policy $\pi_{\theta}(y\mid x)$, optimization is performed by minimizing $\mathbb{D}_{\text{KL}}[\pi_{\theta}(y|x) \mid \mid \pi^*(y\mid x)]$, where $\pi^*$ is derived as the optimal policy in Eq.~\ref{eq:optimum_model} using the reward function $r_{\phi}(y, x)$. Some straightforward algebra leads us to the following optimization criterion:

\begin{equation}\label{eq:AC}
    \max_{\pi_{\theta}}\mathbb{E}_{\pi_{\theta}(y\mid x)}\bigg[\underbrace{r_{\phi}(x, y) -\beta\log\sum_{y}\pi_\text{ref}(y\mid x)\exp\left(\frac{1}{\beta}r_{\phi}(x, y)\right)}_{f(r_{\phi}, \pi_\text{ref}, \beta)} - \underbrace{\beta\log\frac{\pi_{\theta}(y\mid x)}{\pi_\text{ref}(y\mid x)}}_{\text{KL}}\bigg]
\end{equation}

The objective under consideration mirrors that used in earlier studies 
\citep{ziegler2020finetuning, stiennon2022learning, bai2022training, ouyang2022training}, where the DPO-equivalent reward is adopted for the reward class $r_{\phi}$. Here, the normalization component in $f(r_{\phi}, \pi_\text{ref}, \beta)$ can be viewed as the soft value function for the reference policy $\pi_\text{ref}$. Although this normalization term does not influence the set of optimal policies, omitting it may lead to high-variance policy gradients, thereby destabilizing the training process. It is possible to integrate this normalization factor by introducing a learned value function, though this introduces further optimization difficulties. A different approach employed in preceding work is to normalize rewards relative to a human completion baseline, effectively serving as a one-sample Monte Carlo estimator for the normalization term. By contrast, the DPO reparameterization leads to a reward formulation that obviates the need for any external baselines.

\section{Experiments}
This section presents an empirical investigation into the effectiveness of DPO for preference-based policy optimization. We begin by studying a controlled text-generation benchmark to analyze DPO’s efficiency in balancing reward maximization with the minimization of KL-divergence from the reference policy, in comparison with established preference learning methods like PPO. Subsequently, we extend our analysis to encompass larger models and more challenging RLHF scenarios, including tasks in summarization and dialogue. Our findings suggest that DPO generally achieves performance on par with or surpassing robust baselines such as PPO-based RLHF, even when using minimal hyperparameter tuning, and is also competitive with selecting the highest-reward trajectory from $N$ samples. A detailed description of our experimental procedures follows, with supplementary information provided in Appendix~\ref{app:exp_details}.

\textbf{Tasks.} We assess performance across three distinct open-ended text generation challenges. Throughout all studies, algorithms derive a policy from a collection of preference data, denoted as $\mathcal{D}=\bigl\{x^{(i)}, y_w^{(i)}, y_l^{(i)}\bigr\}_{i=1}^N$. In the \textbf{controlled sentiment generation} task, $x$ represents the beginning of a movie review selected from the IMDb dataset \cite{maas-EtAl:2011:ACL-HLT2011}, and the objective is for the policy to output $y$ exhibiting positive sentiment. To enable a rigorous assessment, preference pairs for generated samples are synthesized using a pre-trained sentiment classifier, ensuring $p(\text{positive}\mid x,y_w)>p(\text{positive}\mid x,y_l)$. For the SFT baseline, we conduct additional training of GPT-2-large until convergence using the IMDB training split (for further specifics, see App~\ref{app:sentiment_details}). In the \textbf{summarization} setting, $x$ constitutes a Reddit forum post, and the policy is required to create a summary $y$ capturing the main ideas. Aligning with established practices, we employ the Reddit TL;DR summarization dataset \citep{volske-etal-2017-tl}, together with a set of human-labeled preferences compiled by \citeauthor{stiennon2022learning}. The SFT model in this case is obtained by fine-tuning on summaries written by humans\footnote{\url{https://huggingface.co/CarperAI/openai_summarize_tldr_sft}} within the TRLX \citep{leandro_von_werra_2023_7790115} RLHF framework. Notably, the human preference dataset was originally created by \citeauthor{stiennon2022learning} from samples provided by a different, yet similarly fine-tuned, SFT model. For \textbf{single-turn dialogue}, $x$ is a user prompt, varying from scientific queries to personal advice. Here, the policy is expected to produce a relevant and supportive reply $y$; the Anthropic Helpful and Harmless dialogue dataset \citep{bai2022training} is utilized, comprising 170k human-assistant conversations. Each exchange concludes with two candidate responses from a large, unspecified language model and a label signifying the response preferred by annotators. As there is no pre-existing SFT model in this context, we generate the SFT model by adapting an open-domain language model using only the favored responses.

\textbf{Evaluation.} Our experimental framework adopts two primary strategies for evaluation. For controlled sentiment generation, we assess each algorithm’s performance with respect to the constrained reward maximization goal by plotting the trade-off frontier between obtained reward and KL-divergence from the reference policy; this is possible due to access to the underlying reward signal via a sentiment classifier. In practical scenarios, however, the exact reward function remains inaccessible. Therefore, we measure the \textit{win rate} of each method compared to a baseline, employing GPT-4 as a stand-in for human judgment when evaluating both summary quality (in summarization) and helpfulness (in single-turn dialogue). For the summarization experiments, baseline comparisons are made against reference summaries from the test set; in the dialogue domain, comparisons use the preferred response indicated in the test data. While literature indicates that language models may serve as more reliable automatic evaluators than conventional metrics \citep{Chen2023ExploringTU}, we further conduct a human annotation study—presented in Sec.~\ref{sec:human-judgments}—to empirically substantiate our reliance on GPT-4 for assessment. Our findings demonstrate that GPT-4’s evaluations closely align with those of humans, often with human-GPT-4 agreement matching or surpassing that observed among human annotators themselves.

\textbf{Methods.} Alongside DPO, we benchmark a range of established methods for training language models in line with human preference data. As straightforward baselines, we test zero-shot prompting using \textbf{GPT-J} \citep{gpt-j} for summarization, and 2-shot prompting via \textbf{Pythia-2.8B} \citep{biderman2023pythia} for the dialogue task. We additionally include both the \textbf{SFT} approach and \textbf{Preferred-FT}, in which models are fine-tuned on human-selected outputs $y_w$; this is based on outputs from the SFT model (for sentiment and summarization tasks) or a standard LM (in the case of single-turn dialogue). Another supervised variant, \textbf{Unlikelihood}~\citep{welleck2019neural}, encourages the model to increase the likelihood of $y_w$ while actively decreasing the likelihood of $y_l$, with an adjustable coefficient $\alpha\in[0,1]$ on the penalization term. Furthermore, we incorporate \textbf{PPO} \citep{schulman2017proximal}, which relies on a reward model trained from preference-labeled data, as well as an oracle variant, \textbf{PPO-GT}, which operates with direct access to ground-truth rewards in controlled sentiment settings. For PPO-GT, two implementations are utilized: an off-the-shelf version \cite{leandro_von_werra_2023_7790115}, and an adapted version that includes reward normalization and refined hyperparameter tuning to enhance efficacy (these adjustments are also applied to PPO with learned rewards). Lastly, we evaluate the \textbf{Best of $N$} approach, wherein $N$ responses are sampled from the SFT model (or from Preferred-FT in the dialogue task), and the one with the highest score—according to a reward model trained on preference data—is chosen. While this method achieves strong performance by dissociating reward estimation from PPO optimization, it proves computationally infeasible for even moderate values of $N$ due to the need to generate $N$ samples for every test input.

\subsection{How well can DPO optimize the RLHF objective?}

Typical RLHF methods employ a KL-regularized reward maximization objective, which simultaneously encourages high reward and prevents the policy from straying excessively from the reference policy. Thus, comparing algorithmic performance requires assessing both the magnitude of the reward and the corresponding KL divergence; obtaining marginally greater reward at the cost of a substantial increase in KL is not always advantageous. In Figure~\ref{fig:frontier-tldr-main}, we illustrate the tradeoff curve between reward and KL for several algorithms on the sentiment task. For each algorithm, we perform several independent runs, each using a distinct hyperparameter value governing policy conservativeness (specifically, target KL $\in\{3,6,9,12\}$ for PPO, $\beta \in \{0.05,0.1,1,5\}$, and $\alpha\in\{0.05,0.1,0.5,1\}$ for unlikelihood, with different random seeds for preferred-FT), comprising a total of 22 experiments. Policies are evaluated every 100 training steps until convergence, with performance measured by the mean reward over a held-out set of prompts using the true reward function, and the mean sequence-level KL\footnote{Defined as the aggregate KL divergence across all timesteps in the sequence.} relative to the reference policy, $\text{KL}\left(\pi\mid \mid \pi_\text{ref}\right)$. Our results reveal that DPO consistently attains a reward-KL frontier that is significantly more efficient than those of the other algorithms: DPO achieves the highest reward while maintaining relatively low KL. This superiority is particularly striking for several reasons. Not only do DPO and PPO target the same underlying objective, but DPO demonstrates marked improvements in efficiency, with its reward-KL tradeoff strictly surpassing PPO's across the board. Moreover, DPO outperforms PPO even when the latter is supplied with ground truth rewards (PPO-GT).

\subsection{Can DPO scale to real preference datasets?}
\label{sec:dpo-real-datasets}
We proceed by assessing how DPO performs during fine-tuning on tasks involving summarization and single-turn dialogue. In the context of summarization, it is known that standard automated metrics like ROUGE may not align closely with human preferences~\citep{stiennon2022learning}, and previous studies have indicated that optimizing LMs via PPO using human feedback can yield more effective summaries. For our evaluation, we generate model outputs on the TL;DR summarization dataset's test split, then measure each method's average win rate by comparing generated completions with reference summaries from the test data. Outputs from all models are drawn using temperatures in the range [0.0, 1.0]; the resulting win rates are depicted in Figure~\ref{fig:frontier-tldr-main} (right). All approaches—DPO, PPO, and Preferred-FT—start from the identical GPT-J SFT checkpoint\footnote{\url{https://huggingface.co/CarperAI/openai_summarize_tldr_sft}} for fine-tuning. Our results show that DPO achieves an approximate win rate of 61\% when sampling with temperature 0.0, outperforming PPO, which attains a win rate near 57\% at its best sampling temperature (also 0.0). Furthermore, DPO secures a greater peak win rate in comparison to the $N$ baseline. It is important to mention that we did not conduct a thorough search over DPO’s $\beta$ hyperparameter, suggesting the method’s capabilities may be underrepresented. In addition, DPO displays considerably higher robustness to changes in sampling temperature than PPO, which can regress to the baseline GPT-J model's performance at elevated temperatures. Notably, Preferred-FT provides little to no advantage over the SFT baseline. A direct human evaluation comparing DPO and PPO, presented in Section~\ref{sec:human-judgments}, further reveals that human judges favored DPO samples at temperature 0.25 over PPO samples at temperature 0 in 58\% of cases.

For single-turn dialogue tasks, our evaluation focuses on those examples within the test set of the Anthropic HH dataset \citep{bai2022training} that comprise a single exchange between a human and an assistant. To assess method performance, GPT-4 judges are used, where the annotated preferred completions from the test set serve as references for determining the win rates of each method. Because there is no universally accepted SFT model tailored for this dataset, we initialize from the pre-trained Pythia-2.8B model and use Preferred-FT to fine-tune a reference model on the selected completions, ensuring outputs align with the data distribution. Subsequently, we apply DPO training. For comparison, we also report results using the highest scoring completion out of 128 Preferred-FT samples (as we observed no further gains beyond 128 candidates for this benchmark—refer to Appendix Figure~\ref{fig:best-of-n}) and include a 2-shot prompted version of the Pythia-2.8B base model. Across all approaches, DPO matches or surpasses the top win rates, given optimal temperature settings for each. Additionally, we assess an RLHF system trained with PPO on the Anthropic HH dataset\footnote{\url{https://huggingface.co/reciprocate/ppo_hh_pythia-6B}}, sourced from a recognized implementation\footnote{\url{https://github.com/CarperAI/trlx/tree/main/examples/hh}}; however, we are unable to identify any prompt or sampling temperature combination that enables this model to outperform the unmodified Pythia-2.8B base. Drawing on both our TL;DR experiments and the alignment in reward objectives across these methods, we treat the Best of 128 baseline as an approximate benchmark for PPO-level performance. Notably, DPO is the only method that is computationally efficient while consistently surpassing the reference completions in the Anthropic HH dataset, and delivers comparable or improved results relative to the much costlier Best of 128 strategy. Moreover, as shown in Figure~\ref{fig:dialogue-main}, DPO achieves peak performance in relatively few training steps.

\subsection{Generalization to a new input distribution}

To further analyze how PPO and DPO adapt to distributional changes, we tested the policies obtained from our Reddit TL;DR summarization experiment on a distinct dataset: news articles from the CNN/DailyMail test split \citep{nallapati-etal-2016-abstractive}. This evaluation used the optimal sampling temperatures identified for TL;DR (0 and 0.25). Table~\ref{tab:ood} shows the outcomes. To assess performance, we calculated GPT-4 win rates by comparing generated summaries to reference summaries, applying the same GPT-4 (C) prompt as in the Reddit TL;DR task, except that “forum post” was replaced by “news article.” On this out-of-distribution task, DPO maintains a clear performance advantage over PPO. These findings suggest that DPO policies exhibit generalization capabilities at least on par with PPO, despite DPO not having access to the additional unlabeled Reddit TL;DR prompts leveraged by PPO.

\subsection{Validating GPT-4 judgments with human judgments}
\label{sec:human-judgments}
We ran a human evaluation to assess whether GPT-4’s judgments are consistent with human preferences, based on TL;DR summarization outcomes and using two different GPT-4 prompts. The \textbf{GPT-4 (S)} (simple) prompt asks which summary best covers the important information in the post, whereas the \textbf{GPT-4 (C)} (concise) prompt also includes a criterion for conciseness, since GPT-4 tends to select longer and more repetitive summaries compared to human raters with the simple prompt. Full prompt texts are available in Appendix~\ref{app:prompts}. We conducted three pairwise comparisons—using the top-performing (DPO, temp. 0.25), lowest-performing (PPO, temp. 1.0), and an intermediate (SFT, temp. 0.25) system—each compared against PPO sampled greedily at its optimal temperature, to sample a range of summary qualities. Analysis shows that, under both prompt types, GPT-4’s choices correlate with human judgments to a similar extent as human raters agree among themselves, indicating GPT-4 is a reliable stand-in for human evaluation (note: due to rater constraints, multiple human judgments are available only for DPO and PPO-1 pairings). In sum, the \textbf{GPT-4 (C)} prompt yields win rates more closely aligned with human results; hence, this prompt is utilized for primary reporting in Section~\ref{sec:dpo-real-datasets}. Further details regarding the human evaluation process, including the online interface used and participating raters, can be found in Appendix~\ref{app:human-study}.

\section{Discussion}
Learning from preferences offers a robust and scalable approach for developing proficient, aligned language models. In this work, we have presented DPO, a straightforward training method that enables language models to learn from preferences without the need for reinforcement learning. Unlike traditional approaches that reformulate the preference learning task to fit into an RL framework—thereby requiring standard RL algorithms—DPO establishes a direct correspondence between language model policies and reward functions. This mapping allows language models to adhere to human preferences using a simple cross-entropy objective, effectively bypassing both reinforcement learning and loss of generality. DPO attains performance on par with, or superior to, prevalent RLHF methods—including those utilizing PPO—while requiring minimal hyperparameter adjustment. Consequently, DPO substantially lowers the entry barrier for training language models informed by human preferences.

\textbf{Limitations \& Future Work.} Several avenues for further research emerge from our findings. One critical question concerns the extent to which the DPO-trained policy generalizes to data outside the training distribution, especially relative to models trained with explicit reward functions. Preliminary experiments indicate that DPO exhibits generalization properties akin to those seen in PPO-based approaches, yet more extensive investigation is warranted. For instance, does self-labeling during DPO training facilitate effective utilization of unlabeled prompts? Another consideration involves understanding the dynamics of reward over-optimization within direct preference optimization and whether the slight decline observed in Figure~\ref{fig:dialogue-main}-right is attributable to this effect. Although our evaluations focus on models containing up to 6B parameters, there is significant potential for examining DPO’s scalability to cutting-edge, substantially larger models. Regarding assessment strategies, we observe that GPT-4-determined win rates are prompt-sensitive; further research could focus on optimizing prompt design to generate high-fidelity evaluations from automated systems. Ultimately, DPO’s potential extends well beyond language model training from human feedback, encompassing its application to training generative models across diverse modalities.

\end{document}